\chapter{Introduction}
\label{ch:intro}

Parkinson's disease is the fastest-growing neurological disorder worldwide, and falls are among its
most damaging consequences \parencite{gbd2024nervous,bloem2021parkinson}. Most people with the disease
fall at least once a year, and a large share fall repeatedly \parencite{allen2013falls}. Many of these
falls are linked to \emph{freezing of gait} (FoG): short episodes in which the feet seem glued to the
floor despite the intention to walk, typically when turning, starting to walk or passing through a
doorway \parencite{nutt2011fog}. A wearable sensor that recognised such episodes, or the falls that
follow them, could trigger cueing or call for help.

Deep learning on inertial signals has made fall detection look easy. On public benchmarks, recurrent
and convolutional networks routinely report sensitivities and specificities above 95\%
\parencite{musci2021online,wang2019survey}. These benchmarks, however, are built from falls simulated by
young, healthy volunteers onto mattresses \parencite{sucerquia2017sisfall,casilari2017datasets}, and
accuracy on real falls of older people is known to be much lower \parencite{bagala2012realfalls}. Patients
with Parkinson's disease add another layer: their normal gait already differs from that of healthy
volunteers, and the events that matter clinically, such as freezing, do not look like the falls the
models were trained on.

This report asks a direct question: does a neural network trained to detect falls on healthy
subjects (SisFall) transfer, without retraining, to the detection of FoG in patients (Daphnet)? The
experiment is easy to set up and easy to misread. File-level fall labels that include walking, a
waist sensor compared against an ankle sensor, an interpolated target domain or a different training
regime for each model would all mix confounders that have nothing to do with Parkinson's disease into
the population and task differences under study. A negative result obtained under such conditions
says little about \emph{why} transfer fails, so the protocol is designed to exclude them.

\section{Problem statement}

We study zero-shot transfer from simulated falls in healthy people to freezing of gait in
Parkinson's disease, with a single trunk-worn accelerometer, under conditions designed so that the
only differences between the domains are the population and the event. Within that setting we ask:
\begin{enumerate}
  \item[\textbf{RQ1}] Under identical training conditions, how do a CNN-1D, an LSTM, a GRU and a patch
  Transformer compare on subject-independent fall detection?
  \item[\textbf{RQ2}] Do these fall detectors rank FoG windows above normal Parkinsonian walking, with no
  patient labels at all?
  \item[\textbf{RQ3}] How much of the transfer gap is covariate shift that label-free normalisation could
  remove, and how much is concept shift?
  \item[\textbf{RQ4}] When a few patients are labelled, does fall pre-training help?
\end{enumerate}

\section{Contributions}

\begin{itemize}
  \item A \textbf{controlled transfer protocol} between SisFall and Daphnet: matched sensor site, units and
  sampling rate; an anatomical frame identified from the data; impact-centred fall labels; removal of
  sensor dropouts that would otherwise count as walking; and identical training conditions
  (seeds, mini-batch order, augmentation draws, loss, optimiser, tuning budget) for four architectures.
  \item A \textbf{negative result with a diagnosis}. All four models reach an AUROC of 0.988--0.999 on unseen
  SisFall subjects. On Daphnet they reach 0.49--0.63 and almost never fire at their validated
  threshold, while false alarms on walking rise 15- to 18-fold. A hand-crafted freeze index ranks FoG
  better than any of them without training.
  \item A \textbf{decomposition of the shift}. Re-normalising the target inputs changes AUROC by at most
  $\pm0.13$, with no consistent gain. A linear probe still separates healthy and Parkinsonian movement
  in every learned representation, and the direction that encodes ``fall'' is nearly orthogonal to the
  one that would encode ``freeze''. Fine-tuning from fall weights is worse than training from scratch
  in six of eight comparisons, and in all four at the larger label budget.
  \item A \textbf{reproducible pipeline}: executed notebooks, checkpoints that store their own scaler and
  threshold, every table generated from result files, and a one-click GPU run.
\end{itemize}

The remainder of the report is organised as follows. \Cref{ch:related} reviews fall and FoG
detection, sequence architectures and the theory of dataset shift. \Cref{ch:method} formalises the
problem and the protocol. \Cref{ch:setup} describes the data and the exploratory analysis.
\Cref{ch:results} presents the experiments, and \Cref{ch:conclusion} concludes.
